% Chapter — Planetary and lunar ephemerides: trimmed DE440 Chebyshev
% evaluation. Follows the mandatory FR-29 six-section template: purpose;
% assumptions and domain of validity; derivation; discretization and
% implementation; validation evidence; references. The equation labels
% eq:ephemeris:timescale, eq:ephemeris:recurrence, and eq:ephemeris:series
% are echoed verbatim in comments in cpp/src/ephemeris.cpp and in
% python/star_reacher/data_fetch.py (FR-29 equation-label-to-code
% traceability); renaming them breaks that trace.
\chapter{Planetary and Lunar Ephemerides}
\label{ch:ephemeris}

\section{Purpose}
\label{sec:ephemeris:purpose}

This chapter documents the ephemeris service (FR-4, decision D-8): the
positions and velocities of the Sun, the Earth--Moon barycenter (EMB), the
Earth, the Moon, and the Venus, Mars, and Jupiter barycenters, plus the
lunar libration angles, over the span 2020--2060, for use by the force
models, the frame chain (Moon principal-axis frame), and mission analysis.
The data are the Chebyshev coefficients of the JPL DE440 development
ephemeris~\cite{park2021}, copied \emph{verbatim} --- never refitted ---
from the JPL-distributed kernels \texttt{de440s.bsp} (planets and Moon) and
\texttt{moon\_pa\_de440\_200625.bpc} (lunar orientation) into the SREPH~v1
container by the \texttt{star data fetch de440s} pipeline
(\texttt{python/star\_reacher/data\_fetch.py}; format specification in
\texttt{docs/formats/sreph\_v1.md}). The in-core evaluator is implemented in
\texttt{cpp/src/ephemeris.cpp}; because the coefficients are verbatim, the
container and evaluator contribute only floating-point evaluation error and
the accuracy of every returned quantity is inherited from DE440 itself.

\section{Assumptions and domain of validity}
\label{sec:ephemeris:domain}

Assumptions:
\begin{enumerate}
  \item \textbf{Time argument.} All epochs are TDB seconds since J2000 TDB
    (JD 2451545.0 TDB), carried as one binary64 value --- the native epoch
    representation of the source kernels. Conversion from the simulator's
    internal timescale (TAI, decision D-6) to TDB is the time-systems
    chapter's responsibility, not this model's.
  \item \textbf{Frames and centers.} Position components are resolved along
    ICRF axes, the inertial axes of DE440~\cite{park2021}; the GCRF of
    Table~\ref{tab:notation:frames} shares these axes. Each body is stored
    relative to its DE440 center: the solar-system barycenter for the Sun,
    EMB, Venus, Mars, and Jupiter entries; the EMB for the Earth and Moon
    entries. The geocentric Moon state is composed at evaluation time as the
    difference of the EMB-relative Moon and Earth states, so no coefficient
    is refitted.
  \item \textbf{Geometric states.} Returned positions are instantaneous
    geometric states with no light-time or aberration correction, per the
    FR-23 light-time policy (dynamics use geometric ephemeris positions;
    optical sensor models apply their own corrections and bounds).
  \item \textbf{Lunar orientation.} The libration output is the DE440 lunar
    principal-axis orientation: the 3-1-3 Euler angles $\phi$, $\theta$,
    $\psi$ of the Moon principal-axis frame relative to the ICRF equator,
    with $\psi$ accumulating without wrapping, exactly as the source kernel
    stores it. Construction of the Moon PA rotation matrix from these angles
    belongs to the frames chapter.
\end{enumerate}

Domain of validity: epochs within the records stored in the loaded file.
The standard repack covers 2020-01-01 through 2060-01-01 TDB, widened
outward to whole Chebyshev-record boundaries per body (records are never
split); the loader accepts user-supplied wider trims (FR-4).

Out-of-domain behavior, matching the code: an epoch outside every stored
record of the queried body throws \texttt{std::out\_of\_range} naming the
body, the epoch, and the stored coverage; an unknown body name throws
\texttt{std::invalid\_argument} listing the available bodies. The evaluator
\emph{never} extrapolates a polynomial outside its record --- Chebyshev
extrapolation diverges rapidly and silently, so the guard is a hard error
rather than a documented error growth.

Two accuracy boundaries inherited from the data are part of this domain
statement. First, DE440 itself: its planetary and lunar trajectories carry
the fit uncertainties documented by Park et~al.~\cite{park2021}; nothing in
this simulator improves on them. Second, comparisons against ephemeris
services that serve DE441 --- JPL Horizons among them --- see the published
DE440/DE441 lunar difference: DE441 omits the lunar core--mantle tidal
damping term so it can span 30\,000 years, and its lunar orbit consequently
differs from DE440 by under \qty{2}{\metre} across 1970--2020, growing to
roughly \qty{10}{\metre} per century away from the present, predominantly
along-track~\cite{park2021}. The measured 1.8--\qty{5.4}{\metre} difference
between this repack and DE441-sourced Horizons lunar states over 2020--2060
(Section~\ref{sec:ephemeris:validation}) is that documented model
difference, not a repack error.

\section{Derivation}
\label{sec:ephemeris:derivation}

\subsection{Chebyshev representation of DE ephemerides}

The JPL development ephemerides are distributed as piecewise Chebyshev
polynomial fits~\cite{newhall1989,park2021}. Each body's trajectory is
divided into contiguous records (granules) of fixed length $L$ seconds; the
DE440 values are $L = 4$~days for the Moon and Earth (EMB-relative), 16~days
for the Sun, EMB, and Venus, 32~days for the Mars and Jupiter barycenters,
and 8~days for the lunar Euler angles. Within record $k$, beginning at epoch
$t_k$, each component $p$ (a position coordinate in km, or an Euler angle in
rad) is represented as
\begin{equation}
  p(t) = \sum_{j=0}^{N-1} c_j\, T_j(x),
  \qquad
  x = \frac{2\,(t - t_k)}{L} - 1 \;\in\; [-1, 1],
  \label{eq:ephemeris:timescale}
\end{equation}
where $T_j$ is the Chebyshev polynomial of the first kind of degree $j$ and
$N$ is the per-segment coefficient count ($N \le 14$ for every DE440
segment). The coefficients $c_j$ are produced by JPL's constrained
least-squares fit to the numerically integrated ephemeris; the fit
constrains position \emph{and} velocity at every record boundary, so the
piecewise representation is continuous and differentiable across
records~\cite{newhall1989}. This constraint is what the record-boundary
continuity tests of Section~\ref{sec:ephemeris:validation} exercise.

Chebyshev polynomials satisfy $T_j(\cos\vartheta) = \cos(j\vartheta)$. With
$x = \cos\vartheta$, the cosine addition identity
$\cos((j{+}1)\vartheta) + \cos((j{-}1)\vartheta)
 = 2\cos\vartheta\,\cos(j\vartheta)$
gives the three-term recurrence
\begin{equation}
  T_0(x) = 1,
  \qquad
  T_1(x) = x,
  \qquad
  T_{j}(x) = 2x\,T_{j-1}(x) - T_{j-2}(x),
  \quad j \ge 2.
  \label{eq:ephemeris:recurrence}
\end{equation}
Differentiating equation~\eqref{eq:ephemeris:recurrence} with respect to $x$
yields the companion recurrence for the derivatives
$T'_j \equiv \mathrm{d}T_j/\mathrm{d}x$:
\begin{equation}
  T'_0(x) = 0,
  \qquad
  T'_1(x) = 1,
  \qquad
  T'_{j}(x) = 2\,T_{j-1}(x) + 2x\,T'_{j-1}(x) - T'_{j-2}(x),
  \quad j \ge 2.
  \label{eq:ephemeris:derivrecurrence}
\end{equation}
Both recurrences are numerically benign on $x \in [-1,1]$, where
$|T_j(x)| \le 1$. The value and time derivative of the component follow by
direct summation and the chain rule
$\mathrm{d}x/\mathrm{d}t = 2/L$ from equation~\eqref{eq:ephemeris:timescale}:
\begin{equation}
  \boxed{\;
  p(t) = \sum_{j=0}^{N-1} c_j\, T_j(x),
  \qquad
  \dot{p}(t) = \frac{2}{L} \sum_{j=0}^{N-1} c_j\, T'_j(x).
  \;}
  \label{eq:ephemeris:series}
\end{equation}
This ascending-recurrence evaluation is the scheme used by JPL's own
ephemeris interpolators for the same data~\cite{newhall1989}.

\subsection{Record selection}

For a query epoch $t$ inside a segment whose records begin at $t_0$ with
interval $L$, the evaluator selects record
\begin{equation}
  k = \left\lfloor \frac{t - t_0}{L} \right\rfloor,
  \qquad
  \text{clamped to } k = n - 1 \text{ when } t \text{ equals the segment end},
  \label{eq:ephemeris:selection}
\end{equation}
where $n$ is the record count: a boundary epoch shared by records $k-1$ and
$k$ belongs to the record that \emph{begins} there, and the final epoch of
the segment evaluates in the last record at $x = +1$ exactly. Because the
fit constrains both position and velocity at boundaries~\cite{newhall1989},
either record yields the same values to floating-point precision; the fixed
rule exists so that every implementation selects the same record and
reproduces bit-identical results (D-10).

\subsection{Units}

SPK coefficients are in km (positions) and PCK coefficients in rad
(angles)~\cite{naifspk}; coefficients are stored verbatim in those units,
because rescaling them would round every coefficient. Kind-0 outputs are
converted to SI at evaluation output by a single multiplication per
component: $r_{\mathrm{m}} = 1000 \cdot r_{\mathrm{km}}$ and
$v_{\mathrm{m/s}} = 1000 \cdot v_{\mathrm{km/s}}$ (the factor 1000 is exactly
representable, so the conversion costs one correctly rounded operation).

\section{Discretization and implementation}
\label{sec:ephemeris:implementation}

\begin{enumerate}
  \item \textbf{Module and traceability.} The evaluator lives in
    \texttt{cpp/src/ephemeris.cpp}; the repack pipeline and the Python
    reference evaluator live in \texttt{python/star\_reacher/data\_fetch.py}.
    Both echo the labels \texttt{eq:ephemeris:timescale},
    \texttt{eq:ephemeris:recurrence}, and \texttt{eq:ephemeris:series}
    verbatim at the corresponding code.
  \item \textbf{Verbatim repack.} The pipeline copies the Type~2 coefficient
    words byte-for-byte from the DAF arrays of the source kernels into the
    SREPH~v1 container, dropping only the redundant per-record \textsc{mid}
    and \textsc{radius} words after verifying they equal
    $t_0 + (k + \tfrac12)L$ and $L/2$ (constant-interval DE
    segments~\cite{naifspk}). Both source kernels are pinned by SHA-256 at
    download, and the digests ride in the SREPH header for auditability.
  \item \textbf{Executable specification and bit-identity.} The Python
    reference evaluator implements
    equations~\eqref{eq:ephemeris:timescale}--\eqref{eq:ephemeris:selection}
    in the identical operation sequence as the C++ evaluator: same record
    selection, same recurrences, same fixed ascending accumulation order for
    the sums of equation~\eqref{eq:ephemeris:series}, same final scalings.
    Under the D-10 build flags (no fast-math, floating-point contraction
    off) every operation is a correctly rounded IEEE-754 binary64 operation
    in both languages, so the two implementations agree \emph{bit for bit};
    the committed golden vectors gate exactly this property, which in turn
    pins the evaluator across platforms and compilers.
  \item \textbf{Determinism.} Evaluation uses fixed-size stack buffers (the
    coefficient count is bounded by 14 in DE440 and validated against a
    hard cap of 32 at load), a single fixed evaluation order, and no
    reductions or library transcendentals; the only \texttt{libm} call is
    \texttt{std::floor}, which is exact.
  \item \textbf{Epoch-quantization budget.} Near the 2060 end of the span,
    $|t| \approx \qty{1.9e9}{\second}$ lies in the binade $[2^{30}, 2^{31})$,
    so one ulp of the binary64 epoch is
    $2^{30-52} \approx \qty{2.4e-7}{\second}$. At the fastest
    stored motion (Venus barycenter, $\approx \qty{35}{\kilo\metre\per\second}$
    about the SSB) an epoch rounded by half an ulp displaces the returned
    position by at most $\approx \qty{8}{\milli\metre}$ --- three orders
    below the \qty{1}{\metre} validation gate. The boundary-continuity
    measurements of Section~\ref{sec:ephemeris:validation} reproduce this
    number: the observed step across one epoch ulp at a record boundary is
    the true motion over that interval.
  \item \textbf{No text parsing in the core.} The SREPH container is a
    fixed-layout little-endian binary file (96-byte header, 64-byte
    directory entries, coefficient blocks); the loader decodes it with
    fixed-width reads and rejects malformed files --- wrong magic,
    unsupported major version, truncated directory, or a coefficient block
    running past end of file --- with specific error messages
    (D-2 boundary rule; format specification in
    \texttt{docs/formats/sreph\_v1.md}).
\end{enumerate}

\section{Validation evidence}
\label{sec:ephemeris:validation}

Table~\ref{tab:ephemeris:validation} lists the tests that constitute this
chapter's validation evidence. Each test identifier is machine-checked to
exist in the test tree by \texttt{scripts/lint\_chapters.py}; tolerances and
their derivations are recorded next to each test and in
\texttt{tests/golden/ephemeris/manifest.toml}. The full-span comparison of
the complete repack against JPL Horizons --- all 21 epochs for every
quantity, executed once maintainer-side --- is committed as
\texttt{tests/golden/ephemeris/full\_span\_validation.md} with the raw
Horizons transcripts; the committed excerpts reproduce the same comparisons
offline in CI. Horizons serves the Moon from DE441, so the lunar rows
carry the published DE440/DE441 difference
(Section~\ref{sec:ephemeris:domain}) and the authoritative sub-meter DE440
lunar gate runs against an independent evaluation of the checksummed source
kernel.

\begin{table}[htbp]
  \centering
  \caption{Validation evidence for the ephemeris evaluator.}
  \label{tab:ephemeris:validation}
  \begin{tabular}{lp{3.2cm}p{6.2cm}}
    \toprule
    Test ID & Suite & Acceptance basis \\
    \midrule
    \testid{ephemeris_bitlevel_golden} &
      doctest (\texttt{cpp/tests/}) &
      Bit-exact agreement (binary64 equality) with the Python reference
      evaluator at 40 golden cases --- span endpoints, mid-record epochs,
      and two record boundaries per stored segment --- every value
      cross-checked against jplephem at generation time. \\
    \testid{ephemeris_segment_boundary_continuity} &
      doctest (\texttt{cpp/tests/}) &
      Evaluation at an interior record boundary versus
      \texttt{nextafter} below it: position step
      $< \qty{5e-2}{\metre}$ (measured \qty{8.4e-3}{\metre}, the true
      motion across one epoch ulp), velocity step
      $< \qty{5e-8}{\metre\per\second}$; librations
      $< \qty{1e-9}{\radian}$, $< \qty{1e-19}{\radian\per\second}$. \\
    \testid{ephemeris_error_paths} &
      doctest (\texttt{cpp/tests/}) &
      Unknown body $\to$ \texttt{std::invalid\_argument} naming it;
      out-of-span and between-excerpt-run epochs $\to$
      \texttt{std::out\_of\_range}; bad magic, truncated directory,
      truncated coefficients, missing file $\to$ specific
      \texttt{std::runtime\_error}. \\
    \testid{test_repacked_positions_match_horizons} &
      pytest (\texttt{tests/python/}) &
      Repacked positions versus committed JPL Horizons geometric ICRF
      vectors at 21 TDB-midnight epochs per quantity spanning 2020--2060,
      including span endpoints and record boundaries:
      $< \qty{1}{\metre}$ for Sun, EMB, Venus, Mars, Jupiter, Earth
      (measured worst \qty{6.6e-2}{\metre});
      $< \qty{10}{\metre}$ DE441-envelope for the lunar quantities
      (measured 1.8--\qty{5.4}{\metre}, the published DE440/DE441 lunar
      difference~\cite{park2021}). \\
    \testid{test_moon_matches_de440_jplephem_reference} &
      pytest (\texttt{tests/python/}) &
      Lunar segments versus jplephem's independent evaluation of the
      checksummed DE440 kernel at the same 21 epochs:
      $< \qty{1e-3}{\metre}$ (measured worst \qty{8.6e-8}{\metre}) --- the
      authoritative DE440 lunar fidelity gate. \\
    \testid{test_lunar_librations_match_jplephem} &
      pytest (\texttt{tests/python/}) &
      Libration angles and rates versus jplephem's evaluation of the DE440
      lunar PCK at 21 epochs: $< \qty{1e-10}{\radian}$ and
      $< \qty{1e-19}{\radian\per\second}$ (measured \qty{9.1e-13}{\radian},
      \qty{8.5e-22}{\radian\per\second}). \\
    \testid{test_ephemeris_error_paths_via_bindings} &
      pytest (\texttt{tests/python/}) &
      The C++ error contract crosses pybind11 as typed Python exceptions
      (\texttt{ValueError}, \texttt{IndexError}, \texttt{RuntimeError}). \\
    \bottomrule
  \end{tabular}
\end{table}

\section{References}
\label{sec:ephemeris:references}

This chapter relies on Park, Folkner, Williams, and
Boggs~\cite{park2021} for DE440/DE441 --- data content, reference frame, and
the lunar tidal-damping difference between the two ephemerides; on
Newhall~\cite{newhall1989} for the Chebyshev representation of JPL
ephemerides, the boundary-constrained fitting, and the interpolation scheme;
on the NAIF SPK documentation~\cite{naifspk} for the Type~2 kernel record
layout and units; on Giorgini et~al.~\cite{giorgini1996} for the Horizons
system used as the external comparison authority; and on the jplephem
package~\cite{jplephem} as the independent evaluation used to cross-check
every committed golden value. Full bibliographic details appear in the
bibliography.
